\section{Local covers and hyperbolic path graphs}\label{sec:models}

The geometric construction gives relations among the loops in the
surface.  To show that these relations leave an element of infinite
order, we will read paths in auxiliary covers of finite arrangements.
This section constructs those covers and proves the two properties
needed for Lemma~\ref{thm:path-lemma}: their kernels agree on sufficiently
buffered overlaps, and their path graphs are uniformly hyperbolic.
All constants in this section depend only on the finite arrangement
$\cD$.  In particular, they are fixed before choosing the arithmetic
level.

\subsection{Dihedral labels on a finite arrangement}

For a subset $T\subseteq\cD$, set
\[
 U_T=\BB\setminus\bigcup_{H\in T}H.
\]
There are characters $\epsilon_z,\epsilon_w:\pi_1(U_T)\to\Z/2$,
with $\epsilon_z$ equal to one on a meridian of a $z$ hyperplane and
zero on a meridian of a $w$ hyperplane, and conversely for
$\epsilon_w$.  One can define each character by the parity of the
winding number of the product of defining affine equations of the
hyperplanes in its family.  Let $E_T\to U_T$ denote the associated
four-sheet cover; it may be disconnected.  These signs restrict
canonically to any ball patch, up to a choice of sheets.  Indeed,
meridians normally generate the complement group in a simply connected
ball, as is seen by filling a loop with a disk transverse to the
hypersurfaces.  Thus their values determine a sign character there.

We refine these signs using the infinite dihedral group
\[
 D_\infty=\langle r,b\mid r^2=b^2=1\rangle,
 \qquad \epsilon(r)=\epsilon(b)=1.
\]
Its even subgroup is the infinite cyclic translation subgroup
$\langle rb\rangle$.  Conjugation in $D_\infty$ preserves this subgroup
and changes its integer coordinate at most by a sign.

We first construct $q_z:\pi_1(U_T)\to D_\infty$.  The construction
ignores all $w$ hyperplanes.  If $T$ contains fewer than all six $z$
slope hyperplanes, define
\[
 q_z(\gamma)=r^{\epsilon_z(\gamma)}.
\]
Suppose instead that the whole $z$ pencil is present.  With only its
six hyperplanes removed, projection to the direction
$\lambda=z_2/z_1$ gives a map to the six-punctured projective line.
Choose its standard meridians so that the four red punctures occur
first, based together in a small red disk, and the two blue punctures
occur last.  Assign $r$ to each red meridian and $b$ to each blue
meridian.  The sphere relation is respected because $r^4b^2=1$.
We fix this choice of the direction-line representation once, and use
the same choice for every complete pencil and for both families.

It remains to add the coupling hyperplanes that belong to $T$.
They lie over the red-direction disk: on $z_1=\pm c$ in the ball,
\[
 |\lambda|<0.7.
\]
Over $|\lambda|<1.5$, the pencil local system has a reduction of
structure group to $\langle r\rangle$.  In this reduction, multiply
its transition functions by the two-sheet systems defined, for the
coupling hyperplanes that are present, by
\begin{equation}\label{mod:coupling-functions}
 (1-z_1/c)^{1/2},\qquad (1+z_1/c)^{1/2}.
\end{equation}
This means adding their winding parities to the exponent of $r$.
Over $|\lambda|>1$ one has $|z_1|<c$, so both systems in
\eqref{mod:coupling-functions} have the distinguished trivializations
obtained by the branch of the square root near $1$.  These
trivializations identify the modified system with the unmodified
pencil system on $1<|\lambda|<1.5$.  They therefore glue the two
systems.  This constructs $q_z$, up to conjugacy, on all of $U_T$.
Its parity is $\epsilon_z$: this is true on every meridian, and
meridians normally generate.  Interchanging $z$ and $w$ gives $q_w$.

On each component of $E_T$, both representations take values in
translations.  Choose integer coordinates on the translation groups
and write the resulting homomorphisms as
\[
 \ell_z,\ell_w:\pi_1(E_T)\longrightarrow\Z.
\]
The notation is componentwise.  Their individual kernels are
independent of conjugating the dihedral representations or changing
sheets.  The signs of the integer coordinates are immaterial until
we form a sum in the central model.

For each $T$, apply the intersection-center blowups, parity branching,
and normalization of Section~\ref{geo:fill} to the arrangement $T$ in
$\BB$ and its sign cover $E_T$.  The resulting smooth space is its
\emph{filled finite model}.

\begin{lemma}\label{mod:extension}
The homomorphisms $\ell_z$ and $\ell_w$ extend across the blowups and
branched fillings of the finite model prescribed in
Section~\ref{geo:fill}.
\end{lemma}
\begin{proof}
It suffices to check the meridians of the divisors being added to the
sign cover.  Removing subsets of complex codimension at least two
does not change the fundamental group of a smooth manifold, and
adding a smooth divisor kills its meridian.  These assertions follow
by general position for paths and disks, so they also apply to the
successive local modifications here.

For an incomplete $z$ pencil, $q_z$ is the sign character in
$\langle r\rangle$; hence $\ell_z$ is already zero on the sign cover.
For a complete pencil, a meridian of the exceptional divisor over
$L_z$ circles the common intersection at a fixed direction.  Its
image under direction projection is constant, and the coupling
systems trivialize near $L_z$.  Its $q_z$ label is therefore the
identity.  At an intersection of a coupling hyperplane with a red
slope hyperplane, the two local meridians both have label $r$ in the
same reduction.  The meridian of their exceptional divisor is the
product of these two meridians, and its label is $r^2=1$.  Every
meridian that is filled after branching of order two is the square
of a meridian with reflection label, so again has trivial label.
These are all the complete-pencil centers.  The additional blowups
needed for odd partial pencils cause no difficulty, since the label
on the corresponding sign cover is zero.  The $w$ calculation is
identical, and intersections involving both families are products
of these local calculations.  Thus both labels kill every required
meridian.
\end{proof}

\subsection{The sum cover at the distinguished fiber}

For the full model $T=\cD$, let $\overline E$ denote its filled sign
cover.  Its fiber over $o$ is $F_*=C\times C$, and it has the lift
$s$ of $\sigma(z,w)=(w,-z)$ described in Section~\ref{sec:geometry}.
On $F_*$, the action of $s$ interchanges the factors.  The first
translation label is nonzero on the first factor and zero on the
second; the reverse holds for the second label.  For example, the
product of one red and one blue meridian has even parity and label
$rb$, so its lift to $C$ gives a nonzero translation.  This calculation
on $F_*$ agrees with that on the complement: push a representative
path slightly in a nonzero normal direction, keeping its directions
away from the branch points.  The coupling square roots are
trivial there.

To pass from this filled model to the resolved distinguished fiber,
we need an ordinary covering across the fixed points of the coarse
quotient.  The descent argument below will achieve this by making the
lifted symmetry commute with deck translations.  We will orient the
two labels so that their sum is $s$-invariant.

The construction of $q_z$ is unchanged under $z\mapsto-z$, with
the two coupling hyperplanes interchanged.  Indeed this map fixes
the direction $\lambda$, interchanges the functions in
\eqref{mod:coupling-functions}, and preserves their distinguished
trivializations.  The analogous statement holds for $w$.
Consequently $s$ interchanges the two translation homomorphisms up
to signs.  Orient them by the preceding comparison on $F_*$ so that
\[
 \ell_z\circ s_* =\ell_w,
 \qquad \ell_w\circ s_* =\ell_z.
\]
The nonzero restrictions to the two factors determine these signs.
In particular the homomorphism
\begin{equation}\label{mod:sum-label}
 \ell=\ell_z+\ell_w:\pi_1(\overline E)\longrightarrow\Z
\end{equation}
is $s$-invariant.

\begin{lemma}\label{mod:central}
The regular covering associated with $\ker\ell$ admits an order-four
lift of $s$ commuting with its deck translations.  Its quotient by
this lift is an ordinary topological covering of the coarse space
$\overline E/\langle s\rangle$.  Consequently it pulls back to a
covering of the resolution used over the distinguished fiber.
\end{lemma}
\begin{proof}
Choose an $s$-fixed point of $F_*$, for example a point on its
diagonal, and a point above it in the $\ker\ell$ cover.  Invariance
of $\ell$ gives a unique lift $\widehat s$ fixing this point.  Since
$s^4=1$, its fourth power is a deck transformation fixing a point,
hence is the identity.  Its order is four because it projects to
$s$.  Invariance of the integer label also says that conjugation
by $\widehat s$ acts trivially on the deck group.

For descent to a topological cover, consider a point fixed by a
subgroup $H\subseteq\langle s\rangle$.  On the fiber over that
point, the lifted $H$ action commutes with deck translations, and
therefore acts by translations of this integer torsor.  Each such
translation has finite order because the lift is an action of the
finite group $\langle s\rangle$.  It is thus the identity.  Choose
a sufficiently small invariant neighborhood at the point on which
the original cover is trivial.  The stabilizer acts trivially on
the discrete factor, so its quotient is again a product with that
factor.  Translating these neighborhoods under the finite group
proves that the quotient map is a covering.  Pullback of a covering
along the resolution map remains a covering.
\end{proof}

The model using both kernels will be called an \emph{individual
model}.  The full model using the sum kernel and then the lifted
$s$ quotient will be called the \emph{central model}.  The distinction
is essential: the sum is used only near a distinguished fiber.

\subsection{How kernels compare on bounded balls}

Before constructing graph metrics, we prove that these local labels
can be compared without a global coloring of the arithmetic
arrangement.  There are two possible sources of disagreement:
normal-matching isometries may not be unique, and a larger chart
may contain hyperplanes absent from a smaller one.

First suppose the prescribed normal vectors of a complete pencil
have been matched as in Proposition~\ref{prop:arrangement-local}.
Two distinct slope normals span its positive
normal plane and specify its two coordinate functionals, hence the
ratio $z_2/z_1$.  Thus the direction-line local system is independent
of the choice of an isometry extending the normal match.  If a
coupling normal is also matched, its homogeneous linear equation,
together with those two functionals, specifies the ratio $z_1/c$.
Explicitly, if $X_1$ is the first coordinate functional, $T_0$ is
the last homogeneous coordinate, and $F_c=X_1-cT_0$ defines the
positive coupling hyperplane, then
\[
 z_1/c=\frac{X_1}{cT_0}=\frac{X_1}{X_1-F_c}.
\]
The negative coupling gives the same conclusion with the signs reversed.
The gluing rule \eqref{mod:coupling-functions} is therefore intrinsic
to these matched data as well.  The corresponding assertions hold
for the $w$ family.  Under $\sigma$, the two families are interchanged
and the $w$ coordinates are negated.  If the chosen representative
of a relabeled normal differs by a scalar, use that same scalar on
both sides of the match.  The resulting coordinate ratios transform
as just described, so the individual kernels are carried to the
corresponding individual kernels.  This is the equivariance needed
below; it does not require a globally chosen orientation of a
translation label.

Next, consider a coupling hyperplane absent from a metric ball $Q$.
A metric ball in the ball model is a Euclidean ellipsoid and is
convex.  Its $z_1$ image is therefore convex and avoids the relevant
value $c$ or $-c$.  On that image the corresponding function in
\eqref{mod:coupling-functions} has a single-valued square root.
Where the image meets $|z_1|<c$, choose its sign to agree with the
distinguished square root.  The intersection is convex, so one
choice works throughout it.  If the intersection is empty, there
is no overlap requiring a prescribed sign.  This trivializes the
extra coupling system compatibly with the red-disk gluing.  Hence
including or omitting a coupling hyperplane missing $Q$ gives
isomorphic label systems on $Q$.

\begin{lemma}\label{mod:buffer-bound}
For every $R>0$ there is $B_1(R)$ with the following property.
In any finite model with a complete $z$ pencil, if a loop in the
sign cover projects into $B(x,R)$ and has nonzero $z$ translation,
then
\[
 \dist(x,L_z)\le B_1(R).
\]
The analogous assertion holds for $w$.  The same bounds hold if
the loops must avoid any additional analytic subsets.
\end{lemma}
\begin{proof}
If the ball meets at most one hyperplane of the relevant family,
its complement for that family has trivial or cyclic fundamental
group generated by the one meridian.  One way to see the latter
assertion is to use the convexity of the ball and a transverse
complex coordinate for the hyperplane: its projection has an
ellipse as image, and the centers of the convex fibers give a
continuous section.  The complement thus has the homotopy type
of a punctured ellipse.  Thus its dihedral image lies in a
conjugate of an order-two subgroup.  It has zero translation on
the sign lift.  Omitting hyperplanes from the other family does
not alter this conclusion, because the representation factors
through the complement of the relevant family.

If two distinct slope hyperplanes meet $B(x,R)$, take a unit
negative homogeneous vector representing $x$.  Its inner products
with their fixed positive normals are bounded in terms of $R$:
for a normal $n$ the normalized absolute inner product is
$\sinh\dist(x,H_n)$.  The two normals form a basis of the positive
normal plane of $L_z$.  Their inner-product bounds therefore bound
the positive projection of $x$ onto that plane, and hence bound
$\dist(x,L_z)$.  There are only finitely many pairs of normals,
so this bound is uniform.

For the remaining case, suppose toward a contradiction that there
are balls supporting nonzero translation whose centers leave every
bounded neighborhood of $L_z$.  After passing to a subsequence,
at least one fixed coupling hyperplane meets each ball: otherwise
the preceding two cases apply.  The centers tend to a boundary
point $\xi$, and a fixed-radius hyperbolic ball converges to the
same boundary point as its center.  Thus $\xi$ belongs to the
closure of that coupling hyperplane.  The closures of the two
coupling hyperplanes are disjoint.  At $\xi$ the direction coordinate
is defined and lies strictly inside the red disk, since $z_1=\pm c$
and the boundary estimate is uniform:
\[
 |z_2/z_1|\le\frac{\sqrt{a-c^2}}{c}
 <\frac{\sqrt{1-0.85^2}}{0.85}<0.62<0.7.
\]
Eventually the whole ball lies over this red disk and misses the
other coupling hyperplane.  Its label system consequently reduces
to $\langle r\rangle$, contradicting the assumed nonzero
translation.  This proves the bound.  Restricting the class of
allowed loops does not weaken it.
\end{proof}

\begin{proposition}[Compatibility of kernels]\label{prop:kernel-compatibility}
For every $R>0$ there are constants $B(R)>R$ and $J_0(R)$, depending
only on $\cD$, with the following properties.
\begin{enumerate}
\item Suppose two finite-model descriptions, with their prescribed
normal matches, agree on all hyperplanes meeting $B(x,B(R))$.
Identify their sign covers there by a fixed choice of sheets.
On loops in this sign cover that project into $B(x,R)$, their
individual $z$ kernels agree and their individual $w$ kernels
agree.  More precisely, for each family the corresponding integer
homomorphisms agree up to sign.
The descriptions may include additional hyperplanes outside the
buffered ball.
\item In the full model, if $\dist(x,o)>J_0(R)$, then on such loops
\[
 \ker(\ell_z+\ell_w)=\ker\ell_z\cap\ker\ell_w.
\]
\end{enumerate}
Both statements remain valid after restricting the loops to avoid
any further analytic subsets.  The first statement is equivariant
under the symmetry interchanging the families.
\end{proposition}
\begin{proof}
Choose $B(R)>R+B_1(R)+1$, enlarging it if needed to include the
corresponding bound for both families.  Use the subset of hyperplanes
visible in the buffered ball as an intermediate description.
If a larger description has an incomplete pencil, so does this
subset, and both corresponding translation labels vanish.
If the larger pencil is complete but the buffered pencil is not,
a nonzero translation on an $R$-ball would, by
Lemma~\ref{mod:buffer-bound}, put its center within $B_1(R)$ of the
common pencil intersection.  Since every slope hyperplane contains
that intersection, all six slopes would then meet the buffered
ball, a contradiction.  Thus again both labels vanish on these
loops.  If both pencils are complete, the normal match fixes the
direction-line system, and extra coupling hyperplanes can be
omitted by the square-root trivialization just proved.  The only
remaining change is conjugation of a dihedral label, which changes
a translation coordinate at most by sign.  This proves the first
assertion, including its equivariance and its compatibility with
the chosen sign sheets.

The sets
\[
 \{x:\dist(x,L_z)\le B_1(R)\},\qquad
 \{x:\dist(x,L_w)\le B_1(R)\}
\]
have bounded intersection.  Indeed, in homogeneous coordinates
the two normal planes together form the positive four-dimensional
space.  Bounds on both positive projections bound the distance
from $o$.  Choose $J_0(R)$ beyond that intersection.  Outside it,
at least one of $\ell_z,\ell_w$ vanishes on every loop over
$B(x,R)$, by Lemma~\ref{mod:buffer-bound}.  The displayed equality
of kernels follows.  Every argument concerns the same individual
loops before and after imposing additional omissions, so those
omissions preserve the conclusions.
\end{proof}

\subsection{Graphs defined by projected diameter}

We now turn the local covers into graphs to which the path lemma
applies.  Over each connected component of $E_T$, let
$\widehat E_T$ be the connected regular cover defined by
$\ker\ell_z\cap\ker\ell_w$.  For the full arrangement we also use
the cover defined by $\ker(\ell_z+\ell_w)$ and its lifted-$s$
quotient.  Restrict all these covers to the open arrangement
complement when defining graph vertices.

For $u>0$, vertices are all points of the relevant covering space.
An oriented edge is a parametrized continuous path whose projection
to $\BB$ has diameter strictly less than $u$; its reverse is the
reverse path.  Multiple edges are retained.  In the central case
we take the quotient graph by $\widehat s$.  The symmetry acts
freely on the open arrangement complement, so edges out of a
quotient vertex correspond bijectively to edges out of any chosen
representative.  Denote these graphs by $\mathcal G_T(u)$, with
$\mathcal G_{\mathrm{cen}}(u)$ for the central quotient.

The use of diameter has an important consequence.  A word of
arbitrary length in loops supported in one fixed bounded region
can be one edge.  It is this feature that prevents the unbounded
translation labels from producing large flat regions in the
graphs.

\begin{proposition}[Hyperbolicity of the models]\label{prop:model-hyperbolicity}
There are $u>0$ and $\delta\ge0$, depending only on $\cD$, such that
every connected component of every $\mathcal G_T(u)$ and of
$\mathcal G_{\mathrm{cen}}(u)$ is $\delta$-hyperbolic.
The same $\delta$ can be used if an arbitrary locally finite union
of proper complex analytic subsets of $\BB$ is additionally
omitted, with invariant omissions in the central case.  On the
vertices retained after such an omission, graph distances are
unchanged.
\end{proposition}

The proof has three parts.  First, paths carrying arbitrary deck
translations inside a fixed bounded region give bounded graph diameter
over every bounded radial region.  Next, an arrangement-preserving
collar identifies the remaining vertices with points of a covering of
the sphere complement, together with a radial coordinate; the angular
diameter permitted by a graph edge decays exponentially with that
coordinate.  Finally, comparison with the cone graph defined below
proves uniform hyperbolicity.  No properness or local finiteness of the
graphs is required.

\begin{lemma}\label{mod:bounded-core}
One can choose $u$, simultaneously for all the finite models, so
that the vertices projecting into any fixed bounded radial region
have bounded graph diameter within each connected component.
\end{lemma}
\begin{proof}
The deck group over a component of $E_T$ is a subgroup of $\Z^2$,
or of $\Z$ for the sum cover, and is finitely generated.  At a
basepoint choose loops representing a finite set of its generators.
There are only finitely many sign sheets; also choose finitely
many paths joining the sheets that belong to the same component.
All these projected paths lie in a fixed bounded region.  Choose
$u$ greater than its diameter, for each of the finitely many models.
An arbitrary word in the generator loops then defines a single
edge.  Thus all deck translates over a basepoint have uniformly
bounded graph distance, including the finitely many sign choices.

For a fixed radial bound, any point in the ball region can be
joined to a chosen basepoint by a uniformly bounded chain of small
balls.  Within each ball, paths can avoid the finite complex
hypersurface arrangement: its complement is path connected by
general position.  Choose successive junctions off the arrangement
in the overlaps.  The resulting paths have uniformly bounded
edge count and lift to the covering.  Any discrepancy in their
terminal sheet is corrected using the preceding generator loops.
This gives the asserted diameter bound.  One can take the chain
bound uniformly, for example by covering the compact closure of
the radial region and a path to the basepoint by finitely many
small balls with connected overlap graph.  No positive lower
bound on distance to the arrangement is needed.
\end{proof}

\subsection{The arrangement at the sphere at infinity}

Rescale the ball to unit radius and write
\[
 x=\tanh(r)\xi,\qquad |\xi|=1.
\]
Let $\Sigma$ be its unit sphere.  Each affine intersection of the
hyperplanes that reaches $\Sigma$ is transverse to $\Sigma$.
Indeed its homogeneous Hermitian subspace is nondegenerate by the
arithmetic normal-space property in the proof of
Proposition~\ref{prop:arrangement-local}.  More explicitly, write
an affine intersection as $a_0+V_0$ with $a_0\perp V_0$.  On its
homogeneous span, the Hermitian form is
\[
 |v|^2+(|a_0|^2-1)|t|^2.
\]
Nondegeneracy excludes $|a_0|=1$, precisely the tangency case.
If the intersection reaches the sphere, then $|a_0|<1$, and the
intersection is transverse there.  This also excludes intersections
supported only on the boundary.

There is an arrangement-preserving product collar near $\Sigma$.
Here is a direct construction, including the estimate needed for
the graph metric.  At each $\xi\in\Sigma$, choose a constant real
vector tangent to all affine hyperplanes through $\xi$ and with
positive outward radial component.  Transversality of their
intersection supplies this vector.  Use a small neighborhood
missing every other hyperplane.  A partition of unity combines
these local vectors into a smooth vector field tangent to each
hyperplane and transverse to the spheres.  Normalize it so that
the derivative of the Euclidean radius is one.  Its flow yields,
for $r\ge r_0$ with $r_0$ sufficiently large, diffeomorphisms
between the ideal arrangement complement and the complement on
the radius-$r$ sphere.  They preserve all arrangement strata.
Because the vector field is bounded on a compact collar, their
angular displacement from the original ideal point is
\begin{equation}\label{mod:collar-shift}
 O(1-\tanh r)=O(e^{-2r})
 \quad\hbox{in Euclidean distance.}
\end{equation}
In the central model, average the field under $\sigma$ before
normalizing.  This retains tangency and positive radial component
and makes the collar equivariant.  The collar lifts to each
covering; the lifted end is a radial product with a possibly
disconnected covering $I$ of the ideal complement.

On $\Sigma$ use the metric
\[
 d_K(\xi,\eta)=|1-\langle\xi,\eta\rangle|^{1/2}.
\]
For completeness, it satisfies the triangle inequality.  If
$a=d_K(\xi,\zeta)$ and $b=d_K(\zeta,\eta)$, then
$|\xi-\zeta|\le\sqrt2a$ and
$|\zeta-\eta|\le\sqrt2b$.  Expanding the Hermitian inner product
around $\zeta$ gives
\[
 |1-\langle\xi,\eta\rangle|
 \le a^2+b^2+|\xi-\zeta|\,|\eta-\zeta|
 \le(a+b)^2.
\]
Positivity and symmetry are immediate.  In particular,
\eqref{mod:collar-shift} is an $O(e^{-r})$ displacement in $d_K$.

The ball distance in our normalization satisfies
\begin{equation}\label{mod:ball-distance}
 \cosh d(x,y)=
 \frac{|1-\langle x,y\rangle|}
 {\sqrt{1-|x|^2}\sqrt{1-|y|^2}}.
\end{equation}
Consequently, for points of radii $r+O(1)$, bounded hyperbolic
distance is equivalent, with uniform constants, to angular
$d_K$ distance $O(e^{-r})$.  To check both implications, write
$x=\rho\xi$, $y=\tau\eta$, where $\rho=\tanh r$ and
$\tau=\tanh r'$.  Then
\[
 1-\langle x,y\rangle
 =1-\rho\tau+\rho\tau(1-\langle\xi,\eta\rangle).
\]
The second parenthesis has nonnegative real part.  Thus the
absolute value is at least $\rho\tau d_K(\xi,\eta)^2$ and at most
$1-\rho\tau+d_K(\xi,\eta)^2$.  For $|r-r'|$ bounded and both
radii large, the denominator in \eqref{mod:ball-distance} and
$1-\rho\tau$ are comparable to $e^{-2r}$.  These inequalities
prove the assertion.  Conversely a bounded ball distance bounds
$|r-r'|$ by the triangle inequality for radial distance, so no
separate radial hypothesis is needed in that direction.

On the lifted ideal complement $I$, define
\begin{equation}\label{mod:ideal-metric}
 d_I(\xi',\eta')=
 \min\left\{1,\inf_\alpha\diam_{d_K}(\pi\alpha)\right\},
\end{equation}
where $\alpha$ runs over paths in $I$ joining $\xi'$ to $\eta'$
and $\pi$ is projection to $\Sigma$.  Put $d_I=1$ for points in
different components.  The triangle inequality follows by
concatenation, because the projected paths meet and the diameter
of their union is at most the sum of their diameters.  Distinct
projected endpoints give positivity directly.  For distinct lifts
of the same endpoint, take a small evenly covered neighborhood
in the ideal complement.  A path changing the lift must leave
this neighborhood, so its projected diameter has a positive
lower bound.  Therefore $d_I$ is a metric, bounded by one.

\subsection{A cone over an arbitrary bounded metric space}

We isolate the metric calculation, a discrete form of the familiar
hyperbolic-cone construction; compare \cite[Section~7]{BonkSchramm2000}.
We give the proof for arbitrary bounded metric spaces, as no
compactness hypothesis is available here.  For a metric space $(I,d_I)$
with $d_I\le1$, let $\mathcal C(I)$ have vertices $(\xi',t)$,
$\xi'\in I$, $t\in\Z_{\ge0}$.  Add vertical unit edges between
successive heights at the same point, and horizontal unit edges
at height $t$ whenever $d_I(\xi',\eta')\le e^{-t}$.
In particular all vertices at height zero are pairwise adjacent.

\begin{lemma}\label{mod:cone}
The graph $\mathcal C(I)$ is hyperbolic with an absolute constant,
independent of the bounded metric space $I$.  More precisely, put
\[
 j=\min\{t,t',-\log d_I(\xi',\eta')\},
 \qquad -\log0=+\infty.
\]
Its distance $D$ satisfies
\begin{equation}\label{mod:cone-distance}
 t+t'-2j\le
 D\bigl((\xi',t),(\eta',t')\bigr)
 \le t+t'-2j+3.
\end{equation}
\end{lemma}
\begin{proof}
If $\xi'=\eta'$, both the formula without its additive error and
the claim are immediate from vertical distance.  Otherwise,
descend to height $\lfloor j\rfloor$, take one horizontal edge,
and ascend.  This gives the upper bound.

For the lower bound, let any connecting path have minimum height
$m$ and $n\ge1$ horizontal edges.  The number of its vertical
edges is at least $t+t'-2m$.  By the triangle inequality in $I$,
\[
 d_I(\xi',\eta')\le n e^{-m}.
\]
Since also $m\le t,t'$, this implies $m\le j+\log n$.
The path length is therefore at least
\[
 t+t'-2j+n-2\log n\ge t+t'-2j.
\]
The last inequality holds for every integer $n\ge1$ (indeed its
minimum over positive real $n$ is $2-2\log2>0$).

To deduce hyperbolicity, fix a height-zero basepoint $p$.
Its distance to $(\xi',t)$ differs from $t$ by at most one.
It follows from \eqref{mod:cone-distance} that the Gromov product
of two vertices differs by at most $3/2$ from their corresponding
$j$ value.  For three base points, the metric triangle inequality
implies
\[
 -\log d_I(\xi',\eta')
 \ge\min\{-\log d_I(\xi',\zeta'),
                -\log d_I(\zeta',\eta')\}-\log2.
\]
Taking the minimum with the relevant heights preserves this
inequality.  Thus the Gromov products satisfy the hyperbolicity
inequality with an absolute additive constant.  The standard
Gromov-product characterization of hyperbolicity for graphs
completes the proof; see \cite[Chapter~III.H]{BridsonHaefliger}.
\end{proof}

\subsection{Comparison with the model graphs}

We can now complete the proof of
Proposition~\ref{prop:model-hyperbolicity}.  Work first before the
finite symmetry quotient, in a connected component.  Using the
collar coordinates, send a cone vertex $(\xi',t)$ to the covering
point of radius $r_0+t$ above its ideal coordinate.  A horizontal
edge at positive height is represented by a path whose projected
ideal diameter is at most $2e^{-t}$: the factor two allows the
infimum in \eqref{mod:ideal-metric} to be unattained.  Place that
path at radius $r_0+t$ with the collar.  The angular estimates and
\eqref{mod:ball-distance} give a uniform bound on its projected
ball diameter.  Vertical edges likewise have uniformly bounded
projected diameter.  Enlarge $u$ to exceed these bounds strictly.
At height zero, even when ideal components are different,
Lemma~\ref{mod:bounded-core} gives a uniform bound on the graph
distance of the two images.  The cone map therefore sends adjacent
vertices to pairs at uniformly bounded graph distance.

In the opposite direction, round a point of radius at least $r_0$
to the nearest integer layer, retaining its lifted ideal coordinate.
Map the bounded radial part to one fixed height-zero vertex.
The two maps are coarse inverses.  Radial rounding has uniformly
bounded graph displacement, and the bounded-core lemma handles
all other points.  To verify that this inverse also sends edges
to pairs of bounded distance, consider an edge whose projection
lies sufficiently far out.  If one endpoint has radius $r$, its
whole projected path has radii in $[r-u,r+u]$.  Formula
\eqref{mod:ball-distance} bounds the angular diameter of this path
by $O_u(e^{-r})$.  The collar changes this by another term of the
same order.  Following the lifted path in the collar consequently
gives
\[
 d_I(\xi',\eta')\le C_u e^{-r}.
\]
Its rounded heights differ by a bounded amount, so
\eqref{mod:cone-distance} bounds their cone distance uniformly.
If an edge meets the bounded radial part, both endpoints have
bounded radius and the same conclusion follows through height
zero.  Hence the two graphs are quasi-isometric.  Hyperbolicity
of the cone and quasi-isometry invariance for geodesic spaces
\cite[Chapter~III.H]{BridsonHaefliger} prove hyperbolicity of the
model graph.

For the central model, the collar and the comparison maps on the
end are equivariant under $\widehat s$.  The lifted action preserves
$d_I$, since its projection is unitary and hence preserves $d_K$.
For the finite group $H=\langle\widehat s\rangle$, the formula
\[
 \overline d_I(H\xi',H\eta')
   =\min_{h\in H}d_I(\xi',h\eta')
\]
defines a metric on $I/H$.  Positivity uses finiteness of $H$, and
the triangle inequality follows by translating and concatenating
minimizers.  The quotient of the cone graph is the cone graph of this quotient
metric, possibly with multiple edges and loops, which do not affect
vertex distances.
The preceding comparison thus descends, while the quotient radial
core remains bounded.  Lemma~\ref{mod:cone} proves hyperbolicity
also in the central case.  There are only finitely many subsets
$T$ and one central model, so one choice of $u$ and one
hyperbolicity constant work for all of them.

Finally, let an additional locally finite union of proper complex
analytic subsets be omitted.  Removing vertices and restricting
edges cannot decrease distances.  For the reverse inequality,
take a finite edge path whose endpoints are retained.  Perturb
its finitely many junctions and constituent paths, relative to
the endpoints, to miss the omitted loci.  General position permits
this because proper complex analytic subsets have real codimension
at least two.  Locally finite omissions cause no extra difficulty:
a compact path meets only finitely many members, and the
perturbation can be made in finitely many coordinate neighborhoods.
Make this perturbation as a homotopy of the whole concatenated
chain in the existing arrangement complement, relative to its two
overall endpoints.  Homotopy lifting preserves those two endpoints
in the cover; the intermediate junctions move coherently.  Each
original projected path has diameter strictly smaller than $u$, so a
sufficiently small perturbation preserves that strict inequality.
The perturbed path has exactly the same number of edges.  Thus
all distances between retained vertices are unchanged.  For a
central quotient edge path, lift the whole chain to the cover, perturb it relative to its overall endpoints, and project
back; invariance of the omitted set ensures that the projected
chain avoids it.  The four-point inequality for hyperbolicity
restricts to the retained vertex set with the same constant, proving the
last assertion of the proposition.

We now fix such a value of $u$, a common hyperbolicity constant,
and a reading radius $N$ supplied by Lemma~\ref{thm:path-lemma}.
Only after these choices will we choose the arithmetic level and
the geometric sizes of the charts.
